% SPDX-License-Identifier: BSD-2-Clause
% Copyright (c) 2026 Authors of shirudo
%
% The entry of a kernel function f() in three states: with no program,
% where it holds a nop; with fentry programs only, where the trampoline
% returns into the body; and with an fexit or fmod_ret program, where
% the trampoline calls the body itself and returns to the caller
% (concepts/background/fentry-fexit.rst).
\documentclass[tikz,dvisvgm]{standalone}
\usepackage{shirudo-figure}

\begin{document}
\begin{tikzpicture}
  \useasboundingbox (0,0) rectangle (840,392);

  % No program.
  \draw[box] (10,34) rectangle (130,78);
  \draw[box] (180,34) rectangle (330,78);
  \draw[box] (680,34) rectangle (830,78);
  \draw[flow] (130,56) -- (178,56);
  \draw[flow] (330,56) -- (678,56);

  % fentry programs only.
  \draw[box] (10,138) rectangle (130,182);
  \draw[box] (180,138) rectangle (330,182);
  \draw[box] (380,126) rectangle (630,194);
  \draw[box] (680,138) rectangle (830,182);
  \draw[flow] (130,160) -- (178,160);
  \draw[flow] (330,160) -- (378,160);
  \draw[flow] (630,160) -- (678,160);

  % With an fexit or fmod_ret program.
  \draw[box] (10,262) rectangle (130,306);
  \draw[box] (180,262) rectangle (330,306);
  \draw[box] (380,236) rectangle (630,332);
  \draw[box] (680,262) rectangle (830,306);
  \draw[flow] (130,284) -- (178,284);
  \draw[flow] (330,284) -- (378,284);
  \draw[flow] (630,274) -- (678,274);
  \draw[flow] (680,294) -- (632,294);
  \draw[flow] (505,332) -- (505,352) -- (70,352) -- (70,308);

  \begin{scope}[every node/.append style={small, note, anchor=base west}]
    \node at (10,24) {no program attached};
    \node at (10,116) {fentry programs only};
    \node at (10,226) {with an fexit or fmod\_ret program};
  \end{scope}

  \begin{scope}[every node/.append style={anchor=base}]
    \foreach \y in {61,165,289}
      \node[code] at (70,\y) {f(a, b);};

    \node[title] at (255,52) {entry of f()};
    \node[small, note] at (255,69) {a nop};
    \node[title] at (755,61) {body of f()};

    \node[title] at (255,156) {entry of f()};
    \node[small, note] at (255,173) {calls the trampoline};
    \node[title] at (505,146) {trampoline};
    \node[small] at (505,166) {saves a and b, runs the fentry};
    \node[small] at (505,182) {programs, restores a and b};
    \node[title] at (755,156) {body of f()};
    \node[small, note] at (755,173) {returns to the caller};

    \node[title] at (255,280) {entry of f()};
    \node[small, note] at (255,297) {enters the trampoline};
    \node[title] at (505,256) {trampoline};
    \node[small] at (505,276) {saves a and b, runs the fentry};
    \node[small] at (505,292) {and fmod\_ret programs, calls the};
    \node[small] at (505,308) {body, runs the fexit programs};
    \node[small] at (505,324) {with a, b and the result r};
    \node[title] at (755,280) {body of f()};
    \node[small, note] at (755,297) {returns r to the trampoline};

    \node[small, note] at (288,346) {the trampoline returns r to the caller};
  \end{scope}

  \node[small, note, anchor=base west] at (10,382) {A function has one
    trampoline for all its programs. Each attach and detach generates a
    new image of at most 4 KB.};
\end{tikzpicture}
\end{document}
